\scanpage{36}
\diagram{\begin{tikzcd}[column sep=large,row sep=huge]
t\quad P^1 \arrow[d,"f", "\mathbb Z/3"'] \arrow[d,phantom,"{\deleted{\unsure{section}}}"{description},xshift=18pt] & \underset{\deleted{\unsure{P^1_{\infty}}\ \unsure{G_m}}}{\underbrace{0,1,\infty}} \arrow[d,no head,"{\begin{gathered}\text{non ramifi\unsure{é}}\\\swarrow\qquad\searrow\end{gathered}}"] & \underset{\deleted{\unsure{P^1_m}}}{\underbrace{2,-1,1/2}} \arrow[d,no head] & x'\deleted{\unsure{x}}=\tfrac12(1+\sqrt{-3}) \arrow[d,no head,shift left=3pt] \arrow[d,no head,"{\text{ramifié}}"] \arrow[d,no head,shift right=3pt] & & x''=\tfrac12(1-\sqrt{-3}) \arrow[d,no head,shift left=3pt] \arrow[d,no head,"{\text{ramifié}}"] \arrow[d,no head,shift right=3pt] \\
s\quad P^1 \arrow[d,"g","\mathbb Z/2"'] & \infty \arrow[d,no head,shift left=1.5pt] \arrow[d,no head,shift right=1.5pt,"{\text{ramifié}}"] & 0 \arrow[d,no head,shift left=1.5pt] \arrow[d,no head,shift right=1.5pt,"{\text{ramifié}}"] & 1\deleted{\unsure{x'}} \arrow[dr,no head,"{\deleted{\unsure{\sim}}}"] & \underbrace{\hspace{2.5cm}} \arrow[d,no head,"{\substack{\text{non}\\\text{ramifié}}}"] & -1 \arrow[dl,no head,"{\deleted{\unsure{\sim}}}"] \\
r\quad P^1 & \infty & 0 & & 1 &
\end{tikzcd}}

Pour déterminer explicitement $f$ et $g$, on procède ainsi. [On \deleted{\unsure{se}} peut supposer que \deleted{$S=\operatorname{Spec}(\unsure{\mathbb Z_m}[2^{-1},3^{-1}])$} \inserted{$S=\operatorname{Spec}(A)$, $A=\mathbb Z[2^{-1},3^{-1}]$} \deleted{\unsure{\{}} puisque toute la situation provient par image inverse d'une situation sur $\operatorname{Spec}(A)$]. La restriction de $g$ à $P^1_S-(\infty)-(0)$ est alors un \inserted{$S$-}morphisme de $G_m$ dans lui-même, donc est de la forme $t\to\lambda t^n$, où $n\in\mathbb Z$, et $\lambda\in H^0(S,\underline O_S^*)$. Comme $g(\infty)=\infty$, on a $n>0$, et comme $g$ est \inserted{un rev.} d'ordre 2, on a $n=2$. D'autre part \deleted{$\lambda$} \inserted{$g(1)=1$, d'où $\lambda=1$} \deleted{\unsure{$\in A$, donc de la forme $2^\alpha3^\beta$}.}
\[
(9)\qquad r=g(s)=s^2
\]
\uline{Cette formule est valable sur tout $P^1$} [lui donner un sens, et justifier cette assertion\ldots]
\mspage{35}
